\chapter{Matrices en grafen}\label{ch:g12:matrix}

Een \omterm{def:g12:matrix:matrix}{matrix} is een rechthoekige tabel van getallen, die opgeteld en
vermenigvuldigd worden volgens regels die zo ontworpen zijn dat de algebra van
de matrices het samenstellen van lineaire transformaties weergeeft. Matrices
lossen lineaire stelsels op, sturen gekoppelde recursieve \omterm{def:g12:seq:sequence}{rijen} aan en tellen
wandelingen in netwerken --- de wiskunde achter zoekmachines en algoritmen
voor kortste paden.

\section{Algebra van de matrices}

\begin{definition}[Matrix]\label{def:g12:matrix:matrix}
Een \emph{$m \times n$-matrix}\index{matrix} is een tabel van \omterm{def:g10:numbers:sets}{reële getallen}
met $m$ rijen en $n$ kolommen: $A = (a_{ij})$, waarbij $a_{ij}$ de ingang in
rij $i$, kolom $j$ is. Twee matrices van hetzelfde formaat worden ingang per
ingang opgeteld, en $\lambda A = (\lambda a_{ij})$.
\end{definition}

\begin{definition}[Product van matrices]\label{def:g12:matrix:product}
Zij $A$ een $m \times n$-matrix en $B$ een $n \times p$-matrix. Het product
$AB$ is de $m \times p$-matrix met als ingang $(i,j)$
\[
(AB)_{ij} = \sum_{k=1}^{n} a_{ik} b_{kj}
\]
(de regel ``rij $i$ van $A$ maal kolom $j$ van $B$'').
\end{definition}

\begin{example}
$\begin{pmatrix} 1 & 2\\ 3 & 4\end{pmatrix}
\begin{pmatrix} 0 & 1\\ 1 & 1\end{pmatrix}
= \begin{pmatrix} 2 & 3\\ 4 & 7\end{pmatrix}$,
\quad terwijl \quad
$\begin{pmatrix} 0 & 1\\ 1 & 1\end{pmatrix}
\begin{pmatrix} 1 & 2\\ 3 & 4\end{pmatrix}
= \begin{pmatrix} 3 & 4\\ 4 & 6\end{pmatrix}$:
\emph{het vermenigvuldigen van matrices is niet commutatief}.
\end{example}

\begin{proposition}[Rekenregels voor matrices]\label{prop:g12:matrix:rules}
Telkens wanneer de formaten de producten zinvol maken, geldt
\[
(AB)C = A(BC), \qquad A(B + C) = AB + AC, \qquad (A+B)C = AC + BC,
\]
en de \emph{eenheidsmatrix} $I_n$ (enen op de diagonaal, nullen elders)
voldoet aan $I_m A = A I_n = A$ voor $A$ van formaat $m \times n$.
\end{proposition}

\begin{proof}
Alles zijn verificaties ingang per ingang vanuit
\cref{def:g12:matrix:product}; de associativiteit, als enige niet triviaal,
komt neer op het verwisselen van twee eindige sommen:
$\bigl((AB)C\bigr)_{ij} = \sum_l \left(\sum_k a_{ik}b_{kl}\right) c_{lj}
= \sum_k a_{ik} \left(\sum_l b_{kl} c_{lj}\right)
= \bigl(A(BC)\bigr)_{ij}$.
\end{proof}

\begin{definition}[Inverse]\label{def:g12:matrix:inverse}
Een vierkante \omterm{def:g12:matrix:matrix}{matrix} $A$ van formaat $n$ heet
\emph{inverteerbaar}\index{matrix!inverse} als er een \omterm{def:g12:matrix:matrix}{matrix} $B$ bestaat met
$AB = BA = I_n$; $B$ is dan uniek en wordt $A^{-1}$ genoteerd.
\end{definition}

\begin{proposition}[Inverse van een $2\times2$-matrix]\label{prop:g12:matrix:inverse2x2}
Zij $A = \begin{pmatrix} a & b\\ c & d\end{pmatrix}$ en
$\det A = ad - bc$ (de \emph{determinant}\index{determinant}). Dan is $A$
\omterm{def:g12:matrix:inverse}{inverteerbaar} als en slechts als $\det A \neq 0$, en in dat geval geldt
\[
A^{-1} = \frac{1}{ad - bc}\begin{pmatrix} d & -b\\ -c & a\end{pmatrix}.
\]
\end{proposition}

\begin{proof}
Een berekening geeft
$A \begin{pmatrix} d & -b\\ -c & a\end{pmatrix}
= \begin{pmatrix} d & -b\\ -c & a\end{pmatrix} A = (ad - bc) I_2$; is
$ad - bc \neq 0$, deel dan. Omgekeerd: is $ad - bc = 0$, dan zijn de kolommen
van $A$ evenredig, en dus ook die van $AB$ voor elke $B$; maar de kolommen van
$I_2$ zijn niet evenredig, dus geen enkele $B$ kan aan $AB = I_2$ voldoen.
\end{proof}

\begin{method}[Lineaire stelsels]\label{met:g12:matrix:systems}
Het stelsel $\begin{cases} ax + by = e\\ cx + dy = f \end{cases}$ is de
matrixvergelijking $AX = Y$ met
$X = \begin{pmatrix} x \\ y\end{pmatrix}$ en
$Y = \begin{pmatrix} e \\ f\end{pmatrix}$. Is $\det A \neq 0$, dan is de
unieke oplossing $X = A^{-1}Y$. Hetzelfde formalisme behandelt $n$
\omterm{def:g10:algebra:equation}{vergelijkingen} in $n$ onbekenden.
\end{method}

\section{Machten van matrices en recursieve rijen}

\begin{definition}\label{def:g12:matrix:power}
Voor een vierkante \omterm{def:g12:matrix:matrix}{matrix} $A$ en $k \in \N$ is
$A^k = A \times \dots \times A$ ($k$ factoren), met $A^0 = I$.
\end{definition}

\begin{method}[Diagonaal plus nilpotent, en het diagonaliseerbare geval]\label{met:g12:matrix:powers}
Twee standaardmanieren om $A^k$ te berekenen:
\begin{itemize}
  \item Is $A = \lambda I + N$ met $N^2 = 0$, dan krimpt het binomium van
        Newton (hier geldig, want $I$ en $N$ commuteren) tot twee termen:
        $A^k = \lambda^k I + k \lambda^{k-1} N$.
  \item Vind je een inverteerbare $P$ met $A = PDP^{-1}$ en $D$ diagonaal, dan
        is $A^k = P D^k P^{-1}$, en $D^k$ bereken je ingang per ingang.
        (Zo'n $P$ systematisch vinden is de theorie van het
        \emph{diagonaliseren}, uitgewerkt aan de universiteit; op dit niveau
        wordt $P$ gegeven.)
\end{itemize}
\end{method}

\begin{example}[Gekoppelde rijen]\label{ex:g12:matrix:coupled}
Zij $u_{n+1} = 3u_n + v_n$ en $v_{n+1} = u_n + 3v_n$. Met
$X_n = \begin{pmatrix} u_n\\ v_n \end{pmatrix}$ en
$A = \begin{pmatrix} 3 & 1\\ 1 & 3\end{pmatrix}$ krijgen we
$X_{n+1} = AX_n$, dus $X_n = A^n X_0$. De hulprijen $s_n = u_n + v_n$ en
$d_n = u_n - v_n$ voldoen aan $s_{n+1} = 4s_n$ en $d_{n+1} = 2d_n$, dus is
$s_n = 4^n s_0$ en $d_n = 2^n d_0$, en
\[
u_n = \frac{4^n(u_0+v_0) + 2^n(u_0-v_0)}{2}, \qquad
v_n = \frac{4^n(u_0+v_0) - 2^n(u_0-v_0)}{2}.
\]
(Achter de schermen: $(1,1)$ en $(1,-1)$ zijn de richtingen van de
eigenvectoren van $A$.)
\end{example}

\section{Grafen en wandelingen}

\begin{definition}[Graaf, verbindingsmatrix]\label{def:g12:matrix:graph}
Een \emph{graaf}\index{graaf} bestaat uit toppen $1, 2, \dots, n$ en bogen die
bepaalde paren toppen verbinden (geordende paren voor een \emph{gerichte}
graaf). Zijn \emph{verbindingsmatrix}\index{verbindingsmatrix} is de
$n \times n$-matrix $M$ met $m_{ij} = 1$ als er een boog van $i$ naar $j$
loopt, en $0$ anders. Een \emph{wandeling} van lengte $k$ van $i$ naar $j$ is
een opeenvolging van $k$ bogen die van $i$ naar $j$ leidt.
\end{definition}

\begin{omfigure}
\begin{tikzpicture}[baseline=(current bounding box.center)]
  \node[circle, draw, inner sep=2.5pt] (v1) at (0,0) {$1$};
  \node[circle, draw, inner sep=2.5pt] (v2) at (2.6,1.2) {$2$};
  \node[circle, draw, inner sep=2.5pt] (v3) at (2.6,-1.2) {$3$};
  \draw[-Stealth, omDef, thick] (v1) -- (v2);
  \draw[-Stealth, omDef, thick] (v2) -- (v3);
  \draw[-Stealth, omDef, thick] (v1) to[bend right=18] (v3);
  \draw[-Stealth, omDef, thick] (v3) to[bend right=18] (v1);
\end{tikzpicture}
\qquad\qquad
$M = \begin{pmatrix} 0 & 1 & 1\\ 0 & 0 & 1\\ 1 & 0 & 0 \end{pmatrix}$

{\small Een gerichte \omterm{def:g12:matrix:graph}{graaf} en zijn \omterm{def:g12:matrix:graph}{verbindingsmatrix}
(\cref{exo:g12:matrix:6}): $m_{ij} = 1$ precies wanneer er een boog van $i$
naar $j$ loopt.}
\end{omfigure}

\begin{theorem}[Wandelingen tellen]\label{thm:g12:matrix:walks}
Het aantal wandelingen van lengte $k$ van top $i$ naar top $j$ is de ingang
$(i,j)$ van $M^k$.
\end{theorem}

\begin{proof}
Inductie op $k$. Voor $k = 1$ is dit de definitie van $M$. Onderstel de
bewering voor $k$. Een wandeling van lengte $k+1$ van $i$ naar $j$ is een
wandeling van lengte $k$ van $i$ naar een zekere top $l$, gevolgd door een
boog van $l$ naar $j$; volgens het som- en het productprincipe is hun aantal
\[
\sum_{l=1}^{n} \bigl(M^k\bigr)_{il}\, m_{lj} = \bigl(M^{k+1}\bigr)_{ij}.
\qedhere
\]
\end{proof}

\begin{example}
Voor de driehoeksgraaf ($3$ toppen, alle paren verbonden) is
$M = \begin{pmatrix} 0&1&1\\ 1&0&1\\ 1&1&0\end{pmatrix}$ en
$M^2 = \begin{pmatrix} 2&1&1\\ 1&2&1\\ 1&1&2\end{pmatrix}$: vanuit elke top
zijn er $2$ wandelingen van lengte $2$ terug naar zichzelf (via een van beide
buren) en $1$ naar elke andere top.
\end{example}

\section{Oefeningen}

\begin{exercise}[$\star$]\label{exo:g12:matrix:1}
Zij $A = \begin{pmatrix} 1 & 2\\ 0 & 1 \end{pmatrix}$ en
$B = \begin{pmatrix} 2 & 0\\ 1 & 1 \end{pmatrix}$. Bereken $A + B$, $AB$,
$BA$ en $A^2$.
\end{exercise}

\begin{exercise}[$\star$]\label{exo:g12:matrix:2}
Ga na of de volgende matrices \omterm{def:g12:matrix:inverse}{inverteerbaar} zijn, en bereken de inversen waar
ze bestaan:
\[
A = \begin{pmatrix} 2 & 5\\ 1 & 3\end{pmatrix}, \qquad
B = \begin{pmatrix} 3 & 6\\ 2 & 4\end{pmatrix}.
\]
\end{exercise}

\begin{exercise}[$\star$]\label{exo:g12:matrix:3}
Los het stelsel $\begin{cases} 2x + 5y = 1\\ x + 3y = 2 \end{cases}$ op door
een \omterm{def:g12:matrix:matrix}{matrix} te inverteren.
\end{exercise}

\begin{exercise}[$\star\star$]\label{exo:g12:matrix:4}
Zij $A = \begin{pmatrix} 2 & 1\\ 0 & 2\end{pmatrix} = 2I + N$ met
$N = \begin{pmatrix} 0 & 1\\ 0 & 0 \end{pmatrix}$.
\begin{enumerate}
  \item Ga na dat $N^2 = 0$ en dat $I$ en $N$ commuteren.
  \item Leid $A^k$ af voor alle $k \in \N$ en controleer de formule voor
        $k=2$ met een rechtstreekse berekening.
\end{enumerate}
\end{exercise}

\begin{exercise}[$\star\star$]\label{exo:g12:matrix:5}
Zij $A = \begin{pmatrix} 0 & 1\\ 1 & 1\end{pmatrix}$ en $F_n$ de \omterm{def:g12:seq:sequence}{rij} van
Fibonacci ($F_0 = 0$, $F_1 = 1$, $F_{n+2} = F_{n+1} + F_n$). Toon met inductie
aan dat voor $n \geq 1$ geldt
\[
A^n = \begin{pmatrix} F_{n-1} & F_n\\ F_n & F_{n+1}\end{pmatrix},
\]
en leid de identiteit $F_{n+1}F_{n-1} - F_n^2 = (-1)^n$ af.
(Tip: \omterm{def:g11:vect:det}{determinanten} vermenigvuldigen zich: $\det(MN) = \det M \det N$, wat je
voor $2\times2$-matrices mag nagaan.)
\end{exercise}

\begin{exercise}[$\star\star$]\label{exo:g12:matrix:6}
Een gerichte \omterm{def:g12:matrix:graph}{graaf} op de toppen $\{1, 2, 3\}$ heeft de bogen
$1\to2$, $2\to3$, $3\to1$ en $1\to3$.
\begin{enumerate}
  \item Schrijf de \omterm{def:g12:matrix:graph}{verbindingsmatrix} $M$ op en bereken $M^2$ en $M^3$.
  \item Hoeveel wandelingen van lengte $3$ gaan van $1$ naar $1$? Som ze op.
\end{enumerate}
\end{exercise}

\begin{exercise}[$\star\star$]\label{exo:g12:matrix:7}
Een autodeelbedrijf verplaatst voertuigen tussen twee steden $A$ en $B$. Elke
week blijft $80\%$ van de auto's in $A$ in $A$ en verhuist $20\%$ naar $B$;
van de auto's in $B$ verhuist $30\%$ naar $A$ en blijft $70\%$ ter plaatse.
Zij $a_n$ en $b_n$ de aandelen van het wagenpark in elke stad.
\begin{omfigure}
\begin{tikzpicture}
  \node[circle, draw, inner sep=3pt] (A) at (0,0) {$A$};
  \node[circle, draw, inner sep=3pt] (B) at (3.2,0) {$B$};
  \draw[-Stealth, omDef, thick] (A) to[bend left=25]
    node[above, font=\small] {$0.2$} (B);
  \draw[-Stealth, omDef, thick] (B) to[bend left=25]
    node[below, font=\small] {$0.3$} (A);
  \draw[-Stealth, omDef, thick] (A) to[loop left]
    node[left, font=\small] {$0.8$} (A);
  \draw[-Stealth, omDef, thick] (B) to[loop right]
    node[right, font=\small] {$0.7$} (B);
\end{tikzpicture}
\end{omfigure}
\begin{enumerate}
  \item Schrijf $X_{n+1} = MX_n$ met
        $X_n = \begin{pmatrix} a_n\\ b_n\end{pmatrix}$ en bepaal $M$.
  \item Zoek de evenwichtsaandelen (los $MX = X$ op met $a + b = 1$).
  \item Toon aan dat $c_n = a_n - 0.6$ voldoet aan $c_{n+1} = 0.5\,c_n$, en
        besluit dat de verdeling van het wagenpark naar het evenwicht
        \omterm{def:g12:seq:limit}{convergeert}.
\end{enumerate}
\end{exercise}

\begin{exercise}[$\star\star\star$]\label{exo:g12:matrix:8}
Zij $A = \begin{pmatrix} 3 & 1\\ 1 & 3 \end{pmatrix}$ en
$P = \begin{pmatrix} 1 & 1\\ 1 & -1 \end{pmatrix}$.
\begin{enumerate}
  \item Bereken $P^{-1}$, daarna $D = P^{-1}AP$, en ga na dat $D$ diagonaal
        is.
  \item Leid een gesloten formule voor $A^n$ af en vergelijk met
        \cref{ex:g12:matrix:coupled}.
\end{enumerate}
\end{exercise}

\section{Opgave: de matrix die Fibonacci kent (en het weer)}

\begin{problem}[{Weekendopgave --- één $2 \times 2$-matrix draagt heel
Fibonacci, een matrix van Markov voorspelt het weer op lange termijn, en een
eigenvector is een miljard dollar waard}]
\label{pb:g12:matrix:1}
Een \omterm{def:g12:matrix:matrix}{matrix} is een machine die een toestand opeet en de volgende teruggeeft ---
en haar \emph{machten} bevatten dus hele toekomsten. Deze opgave opent met de
verbluffende \omterm{def:g12:matrix:matrix}{matrix} waarvan de machten de getallen van Fibonacci opsommen (en
hun identiteiten elk in één regel bewijzen), laat daarna het weer als een
keten van Markov naar zijn stationaire toestand lopen, en sluit af met de
eigenvector waarop een zoekmachine gebouwd werd
(\cref{thm:g12:matrix:walks}, \cref{met:g12:matrix:powers}).

\medskip
\noindent\textbf{Deel I --- Vlotheid.}
\begin{enumerate}
  \item Bereken met $A = \begin{pmatrix} 1 & 2\\ 3 & 4\end{pmatrix}$ en
        $B = \begin{pmatrix} 0 & 1\\ 1 & 0\end{pmatrix}$ de producten $AB$ en
        $BA$. Oordeel over de commutativiteit?
  \item Inverteer $\begin{pmatrix} 2 & 1\\ 5 & 3\end{pmatrix}$
        (\cref{prop:g12:matrix:inverse2x2}) en los met die inverse
        $2x + y = 4$, $5x + 3y = 7$ op.
  \item Zij $N = \begin{pmatrix} 0 & 1\\ 0 & 0\end{pmatrix}$: bereken $N^2$ en
        leid $(I + N)^n = I + nN$ af voor elke $n$.
  \item De driehoeksgraaf (drie toppen, alle paren verbonden): schrijf zijn
        \omterm{def:g12:matrix:graph}{verbindingsmatrix} $A$ op, bereken $A^3$, en duid de ingangen op de
        diagonaal (\cref{thm:g12:matrix:walks}).
  \item Geef voor $D = \begin{pmatrix} 2 & 0\\ 0 & \frac12 \end{pmatrix}$ de
        \omterm{def:g12:matrix:matrix}{matrix} $D^n$ en haar gedrag wanneer $n \to \infty$.
\end{enumerate}

\medskip
\noindent\textbf{Deel II --- De \omterm{def:g12:matrix:matrix}{matrix} van Fibonacci.}
Zij $F = \begin{pmatrix} 1 & 1\\ 1 & 0\end{pmatrix}$ en zij
$F_1 = F_2 = 1, F_3 = 2, \dots$ de getallen van Fibonacci uit
\cref{pb:g11:seq:1}.
\begin{enumerate}[resume]
  \item Bereken $F^2$, $F^3$ en $F^4$ en vermoed de algemene vorm van $F^n$ in
        termen van de getallen van Fibonacci.
  \item Bewijs het vermoeden
        $F^n = \begin{pmatrix} F_{n+1} & F_n\\ F_n & F_{n-1} \end{pmatrix}$
        met inductie.
  \item Neem in beide leden de \omterm{def:g11:vect:det}{determinant} (de \omterm{def:g11:vect:det}{determinant} van een product is
        het product van de \omterm{def:g11:vect:det}{determinanten} --- ga het na op
        $2 \times 2$-matrices als je het nooit gezien hebt): leid de
        \emph{identiteit van Cassini} $F_{n+1}F_{n-1} - F_n^2 = (-1)^n$ af ---
        de motor van het verdwijnende vierkantje, in één regel bewezen.
  \item Lees in $F^{m+n} = F^m F^n$ de ingangen rechtsboven af en leid de
        \emph{somformule}
        \[
        F_{m+n} = F_{m+1} F_n + F_m F_{n-1}
        \]
        af. Ga ze na voor $m = n = 3$.
  \item Leid uit de somformule af (inductie op $k$) dat $F_n$ het getal
        $F_{kn}$ \omterm{def:g12:arith:divides}{deelt}, en ga dat na voor $F_3 \mid F_6$ en $F_3 \mid F_9$.
  \item Om $F_{100}$ te berekenen hoef je geen $100$ matrices te
        vermenigvuldigen: kwadrateer herhaaldelijk ($F^2, F^4, F^8, \dots$) en
        combineer. Hoeveel matrixvermenigvuldigingen volstaan er, en welke
        aloude vermenigvuldigingstruc uit het onderbouwvolume is dit, bevorderd
        tot matrices?
\end{enumerate}

\medskip
\noindent\textbf{Deel III --- De weermachine.}
In een zekere stad geldt: na een zonnige dag is de volgende zonnig met \omterm{def:g10:proba:distribution}{kans}
$0.8$; na een regendag is ze zonnig met \omterm{def:g10:proba:distribution}{kans} $0.4$. Codeer de verdeling van
de dag als een kolom $\binom{p_{\text{zon}}}{p_{\text{regen}}}$ en de evolutie
door
\[
M = \begin{pmatrix} 0.8 & 0.4\\ 0.2 & 0.6 \end{pmatrix}.
\]
\begin{enumerate}[resume]
  \item Ga na dat elke kolom van $M$ als som $1$ heeft, en zeg waarom elke
        weermachine die eigenschap moet hebben.
  \item Vandaag is het zonnig. Bereken de voorspelling voor morgen en voor
        overmorgen.
  \item Zoek de \emph{stationaire toestand}: de verdeling $v$ met $Mv = v$
        (en met som $1$). Welk aandeel van de dagen is op lange termijn
        zonnig?
  \item Vertrek van een regendag, $\binom01$, en pas $M$ vier keer toe, waarbij
        je bij elke stap de afstand tot de stationaire toestand bijhoudt. Met
        welke factor krimpt de kloof per stap --- en om welk soort convergentie
        gaat het?
  \item PageRank in het klein: drie pagina's, met de verwijzingen
        $A \to B$, $A \to C$, $B \to C$ en $C \to A$. Een willekeurige surfer
        volgt uniform willekeurig een uitgaande verwijzing. Schrijf de
        overgangsmatrix op, zoek de stationaire toestand, en rangschik de
        pagina's.
  \item Duid de rangschikking: waarom scoort $C$ even hoog als $A$ hoewel het
        van minder pagina's een verwijzing krijgt --- wat meet de stationaire
        toestand eigenlijk? (De echte PageRank voegt een dempingsfactor toe
        voor doodlopende paden en sprongen; het idee van de eigenvector is
        precies dit.)
\end{enumerate}

\medskip
\noindent\textbf{Deel IV --- Het rendement van de diagonaal.}
\begin{enumerate}[resume]
  \item Twee gekoppelde grootheden voldoen aan $u_{n+1} = 3u_n + v_n$ en
        $v_{n+1} = u_n + 3v_n$, dat wil zeggen aan de \omterm{def:g12:matrix:matrix}{matrix} $A$ van
        \cref{exo:g12:matrix:8}. Geef met de diagonalisatie uit die oefening
        ($D = \operatorname{diag}(4, 2)$) de gesloten formule voor $u_n$ als
        $u_0 = 1$ en $v_0 = 0$, en toets ze aan de rechtstreekse berekening
        voor $n = 1, 2, 3$.
  \item In een of twee zinnen: wat \emph{doet} het diagonaliseren met een
        gekoppeld stelsel --- en in welke zin is de stationaire toestand van
        Markov uit vraag 14 ook een verhaal over eigenvectoren?
  \item Slotstuk --- de drie gezichten van de \omterm{def:g12:matrix:matrix}{matrix} dit weekend: boekhouding
        (stelsels en inversen), combinatoriek (wandelingen en verwijzingen
        geteld door machten) en evolutie (Fibonacci, het weer, het web ---
        toekomsten afgelezen op de eigenrichtingen). Telkens één zin, plus de
        blik vooruit: de lineaire algebra van de universitaire volumes maakt
        van elk van die gezichten een theorie.
\end{enumerate}
\end{problem}
